%!TEX root = ./ft-hd.tex
\section{Analysis of \textsc{LocateSignal}: bounding $\ell_1$ norm of undiscovered elements}\label{sec:l1}

The main result of this section is Theorem~\ref{thm:l1-res-loc}, which is our main tool for showing efficiency of \textsc{LocateSignal}. Theorem~\ref{thm:l1-res-loc} applies to the following setting. Fix a set $S\subseteq \nsq$ and a set of hashings $H_1,\ldots, H_{r_{max}}$, and let $S^*\subseteq S$ denote the set of elements of $S$ that are not isolated with respect to most of these hashings $H_1,\ldots, H_{r_{max}}$. Theorem~\ref{thm:l1-res-loc} shows that for any signal $x$ and partially recovered signal $\chi$, if $L$ denotes the output list of an invocation of \textsc{LocateSignal} on the pair $(x, \chi)$ with hashings $H_1,\ldots, H_{r_{max}}$, then 
the $\ell_1$ norm of elements of the residual $(x-\chi)_S$ that are not discovered by \textsc{LocateSignal} can be bounded by a function of the amount of $\ell_1$ mass of the residual that fell outside of the `good' set $S\setminus S^*$, plus the `noise level' $\mu\geq ||x_{\nsq\setminus S}||_\infty$ times $k$.

If we think of applying Theorem~\ref{thm:l1-res-loc} iteratively, we intuitively get that the fixed set of measurements with hashings $\{H_r\}$ allows us to always reduce the $\ell_1$ norm of the residual $x'=x-\chi$ on the `good' set $S\setminus S^*$ to about the amount of mass that is located outside of this good set.


\noindent{\em {\bf Theorem~\ref{thm:l1-res-loc}}
There exist absolute constants $C_1, C_2, C_3>0$ such that for any $x, \chi\in \C^N$ and residual signal $x'=x-\chi$ the following conditions hold. \defsk   Let $S^*\subseteq S$ denote the set of elements that are not isolated with respect to at least a $\sqrt{\alpha}$ fraction of hashings $\{H_r\}_{r=1}^{r_{max}}$.   Suppose that for every $s\in [1:d]$ the sets $\A_r\star(\zero, \mathbf{e}_s)$ are balanced (as per Definition~\ref{def:balance}),  $r=1,\ldots, r_{max}$, and the exponent $F$ of the filter $G$ is even and satisfies $F\geq 2d$.  
Let 
$$
L=\bigcup_{r=1}^{r_{max}}\Call{LocateSignal}{\chi, H_r, \{m(\wh{x}, H_r, a\star (\one, \h)\}_{a\in \A_r, \h\in \H_r}}.
$$


Then if $r_{max}, c_{max}\geq (C_1/\sqrt{\alpha})\log\log N$, one has
$$
||x'_{S\setminus S^*\setminus L}||_1\leq (C_2\alpha)^{d/2} ||x'_S||_1+C_3^{d^2}(||\chi_{\nsq\setminus S}||_1+||x'_{S^*}||_1)+4\mu |S|.
$$
}


As we will show later, Theorem~\ref{thm:l1-res-loc} can be used to show that (assuming perfect estimation) invoking \textsc{LocateSignal} repeatedly allows one to reduce to $\ell_1$ norm of the head elements down to essentially
$$
||x'_{S^*}||_1+||\chi_{\nsq\setminus S}||_1,
$$
i.e. the $\ell_1$ norm of the elements that are not well isolated and the set of new elements created by the process due to false positives in location. 
In what follows we derive bounds on $||e^{head}||_1$ (in section~\ref{sec:h-noise}) and $||e^{tail}||_1$ (in section~\ref{sec:tail-noise}) that lead to a proof of Theorem~\ref{thm:l1-res-loc}.

\subsection{Bounding noise from heavy hitters}\label{sec:h-noise}

We first derive bounds on noise from heavy hitters that a single hashing $H$ results in, i.e. $e^{head}(H, x)$,  (see Lemma~\ref{lm:l1b-single-pi}), and then use these bounds to bound $e^{head}(\{H\}, x)$ (see Lemma~\ref{lm:eh-s-sstar}). These bounds, together with upper bounds on contribution of tail noise from the next section, then lead to a proof of Theorem~\ref{thm:l1-res-loc}.

\begin{lemma}\label{lm:l1b-single-pi}
Let $x, \chi\in \C^N, x'=x-\chi$. \defsk
Let $\pi=(\Sigma, q)$ be a permutation, let $H=(\pi, B, F), F\geq 2d$ be a hashing into $B$ buckets and filter $G$ with sharpness $F$. Let $S^*_H\subseteq S$ denote the set of elements $i\in S$ that are not isolated under $H$. Then one has, for $e^{head}$ defined with respect to $S$, 
$$
||e^{head}_{S\setminus S^*_H}(H, x, \chi)||_1\leq 2^{O(d)} \alpha^{d/2} ||x'_{S\setminus S^*_H}||_1+\GS \cdot 2^{O(d)} (||x'_{S^*}||_1+||\chi_{\nsq\setminus S}||_1).
$$
Furthermore, if $\chi_{\nsq\setminus S}=0$ and $S^*_H=\emptyset$, then one has $||e^{head}_{S}(H, x, \chi)||_\infty\leq 2^{O(d)}\alpha^{d/2} ||x'_S||_\infty$.
\end{lemma}
\begin{proof}
By \eqref{eq:eh-pi} for $i\in S\setminus S^*_H$
\begin{equation}\label{eq:asdf}
\begin{split}
e^{head}_{i}(H, x')&= |G_{o_i(i)}^{-1}|\cdot \sum_{j\in S\setminus S^*_H \setminus  \{i\}} |G_{o_i(j)}|x'_j|\text{~~~~~~~~~~~~~~~~~~~~~~~~~~~~~~~~~~(isolated head elements)}\\
&+|G_{o_i(i)}^{-1}|\cdot \left[\sum_{j\in S^*_H} |G_{o_i(j)}|x'_j|+\sum_{j\in \nsq\setminus S} |G_{o_i(j)}||\chi_j|\right]\text{~~~(non-isolated head elements and false positives)}\\
&=|G_{o_i(i)}^{-1}|\cdot (A_1(i)+A_2(i)).
\end{split}
\end{equation}
Let $A_1:=\sum_{i\in S\setminus S^*_H} A_1(i), A_2:=\sum_{i\in S\setminus S^*_H} A_2(i)$.

We bound $A_1$ and $A_2$ separately. 

\paragraph{Bounding $A_1$.} We start with a convenient upper bound on $A_1$:
\begin{equation}\label{eq:2pogh}
\begin{split}
A_1=&\sum_{i\in S\setminus S^*_H} \sum_{j\in S\setminus S^*_H\setminus  \{i\}} |G_{o_i(j)}||x'_j|~~~~~~~~~~~~~~~~~~~~~~~~~~~~~~~~~~~~~~~~\text{(recall that $o_i(j)=\pi(j)-(n/b)h(i)$)}\\
& =\sum_{t\geq 0} \sum_{i\in S\setminus S^*_H} \sum_{\substack{j\in S\setminus S^*_H\setminus \{i\}\text{~s.t.~}\\ ||\pi(j)-\pi(i)||_\infty\in (n/b)\cdot [2^t-1, 2^{t+1}-1)}} |G_{o_i(j)}||x'_j|, \text{~~~(consider all scales $t\geq 0$)}\\
& \leq  \sum_{t\geq 0} \sum_{i\in S\setminus S^*_H} \max_{||\pi(j)-\pi(i)||_\infty\geq (n/b)\cdot (2^t-1)} G_{o_i(j)}\cdot \sum_{\substack{j\in S\setminus S^*_H\setminus \{i\}\text{~s.t.~}\\ ||\pi(j)-\pi(i)||_\infty\leq (n/b)\cdot (2^{t+1}-1)}} |x'_j|\\
&= \sum_{j\in S\setminus S^*_H} |x'_j|\cdot \sum_{t\geq 0} \max_{\substack{||\pi(j)-\pi(i)||_\infty\geq \\(n/b)\cdot (2^t-1)}} G_{o_i(j)}\cdot \left|\left\{i\in S\setminus S^*_H\setminus \{j\} \text{~s.t.~}||\pi(j)-\pi(i)||_\infty\leq (n/b)\cdot (2^{t+1}-1)\right\}\right|.
\end{split}
\end{equation}
Note that in the first line we summed, over all $i\in S\setminus S^*_H$  (i.e. all isolated $i$), the contributions of all other $i\in S$ to the noise in their buckets. We need to bound the first line in terms of $||x'_{S\setminus S^*_H}||_1$. For that, we first classified all $j\in S\setminus S^*_H$ according to the $\ell_\infty$ distance from $i$ to $j$ (in the second line), then upper bounded the value of the filter $G_{o_i(j)}$ based on the distance $||\pi(i)-\pi(j)||_\infty$, and finally changed order of summation to ensure that the outer summation is a weighted sum of absolute values of $x'_j$ {\em over all $j\in S\setminus S^*_H$}\footnote{We note here that we started by summing over $i$ first and then over $j$, but switched the order of summation to the opposite in the last line. This is because the quantity $G_{o_i(j)}$, which determines contribution of $j\in S$ to the estimation error of $i\in S$ {\em is not symmetric} in $i$ and $j$. Indeed, even though $G$ itself is symmetric around the origin, we have $o_i(j)=\pi(j)-(n/b)h(i)\neq o_j(i)$. }.  In order to upper bound $A_1$ it now suffices to upper bound all factors multiplying $x'_j$ in the last line of the equation above. As we now show, a strong bound follows from isolation properties of $i$. 

We start by upper bounding $G$ using Lemma~\ref{lm:filter-prop}, {\bf (2)}. We first note that by triangle inequality 
$$
||\pi(j)-(n/b)h(i)||_\infty\geq ||\pi(j)-\pi(i)||_\infty-||\pi(i)-(n/b)h(i)||_\infty\geq (n/b)(2^t-1)-(n/b)= (n/b) (2^{t-1}-2).
$$
The rhs is positive for all $t\geq 3$ and for such $t$ satisfies $2^{t-1}-2\leq 2^{t-2}$. We hence get for all  $t\geq 3$
\begin{equation}\label{eq:g-ubound}
\begin{split}
\max_{||\pi(j)-\pi(i)||_\infty\geq (n/b)\cdot (2^{t-1}-1)} G_{o_i(j)}\leq \left(\frac2{1+||\pi(j)-(n/b)h(i)||_\infty}\right)^{\fc}\leq \left(\frac2{1+2^{t-2}}\right)^{\fc}\leq 2^{-(t-3)\fc}.
\end{split}
\end{equation}
We also have the  bound $||G||_\infty\leq 1$ from Lemma~\ref{lm:filter-prop}, {\bf (3)}. It remains to bound the last term on the rhs of the last line in~\eqref{eq:2pogh}. We need the fact that for a pair $i, j$ such that $||\pi(j)-\pi(i)||_\infty\leq 2^{t+1}-1$ we have by triangle inequality
$$
||\pi(j)-(n/b)h(i)||_\infty\leq ||\pi(j)-\pi(i)||_\infty+||\pi(i)-(n/b)h(i)||_\infty\leq (n/b)(2^{t+1}-1)+(n/b)= (n/b) 2^{t+1}.
$$
Equipped with this bound, we now conclude that
\begin{equation}\label{eq:3hrgrg}
\begin{split}
&\left|\left\{i\in S\setminus S^*_H\setminus \{j\} \text{~s.t.~}||\pi(j)-\pi(i)||_\infty\leq (n/b)\cdot (2^{t+1}-1)\right\}\right|\\
&=|\pi(S\setminus \{i\})\cap \B^\infty_{(n/b) h(i)}((n/b)\cdot 2^{t+1})|\leq \GI\cdot \alpha^{d/2} 2^{(t+2)d+1}\cdot 2^t,
\end{split}
\end{equation}
where we used the assumption that $i\in S\setminus S^*_H$ are isolated  (see Definition~\ref{def:isolated}). We thus get for any $j\in S\setminus S^*_H$
\begin{equation*}
\begin{split}
\eta_j:=&\sum_{t\geq 0} \max_{\substack{||\pi(j)-\pi(i)||_\infty\geq \\(n/b)\cdot (2^t-1)}} G_{o_i(j)}\cdot \left|\left\{i\in S\setminus S^*_H\setminus \{j\} \text{~s.t.~}||\pi(j)-\pi(i)||_\infty\leq (n/b)\cdot (2^{t+1}-1)\right\}\right|\\
&\leq \sum_{t\geq 0} (\GI\cdot  \alpha^{d/2} 2^{(t+2)d+1}\cdot 2^{t}) \min\{1,  2^{-(t-3)\fc}\} \\
&\leq \GI\cdot  \alpha^{d/2} 2^{2d+1} \sum_{t\geq 0} 2^{t(d+1)}\cdot \min\{1, 2^{-(t-3)\fc}\}\\
\end{split}
\end{equation*}
We now note that
\begin{equation*}
\begin{split}
\sum_{t\geq 0} 2^{t(d+1)}\cdot \min\{1, 2^{-(t-3)\fc}\}&=1+2^{2(d+1)}+2^{3(d+1)}\sum_{t\geq 3} 2^{(t-3)(d+1)}\cdot \min\{1, 2^{-(t-3)\fc}\}\\
&=1+2^{2(d+1)}+2^{3(d+1)}\sum_{t\geq 3} 2^{(t-3)(d+1-\fc)}\leq 1+2^{2(d+1)}+2^{3(d+1)+1}\leq 2^{4(d+1)+1},\\
\end{split}
\end{equation*}
since $\fc\geq 2d$ by assumption of the lemma, and hence for all $j\in S\setminus S^*_H$ one has $\eta_j\leq  \GI\cdot 2^{O(d)} \alpha^{d/2}$. Combining the estimates above, we now get
\begin{equation*}
\begin{split}
A_1\leq &\sum_{j\in S\setminus S^*_H} |x'_j|\cdot \eta_j\leq  ||x'_S||_1\GI\cdot 2^{O(d)}\alpha^{d/2}, \\
\end{split}
\end{equation*}
 as required.  The $\ell_\infty$ bound for the case when $\chi_{\nsq\setminus S}=0$ follows in a similar manner and is hence omitted.



We now turn to bounding $A_2$. The bound that we get here is weaker since $\chi_{\nsq\setminus S}$ is an adversarially placed signal and we do not have isolation properties with respect to it, resulting in a weaker bound on (the equivalent of) $\eta_j$ for $j\in S^*_H$ than we had for $j\in S\setminus S^*_H$.  We let $y:=x'_{S^*}-\chi_{\nsq\setminus S}$ to simplify notation. We have, as in~\eqref{eq:2pogh},
\begin{equation*}
\begin{split}
A_2&\leq \sum_{j\in S\setminus S^*_H} |x'_j|\cdot \kappa_j,\\
& \text{~where~}\\
&\kappa_j=\sum_{t\geq 0} \max_{\substack{||\pi(j)-\pi(i)||_\infty\geq \\(n/b)\cdot (2^t-1)}} G_{o_i(j)}\cdot \left|\left\{i\in S\setminus S^*_H\setminus \{j\} \text{~s.t.~}||\pi(j)-\pi(i)||_\infty\leq (n/b)\cdot (2^{t+1}-1)\right\}\right|.
\end{split}
\end{equation*}
The first term can be upper bounded as before. For the second term, we note that every pair of points $i_1, i_2\in S\setminus S^*_H$ by triangle inequality satisfy 
$$
(n/b)||\pi(i_1)-\pi(i_2)||_\infty\leq (n/b)||\pi(i_1)-\pi(j)||_\infty+||\pi(j)-\pi(i_2)||_\infty\leq (n/b)\cdot (2^{t+2}-2)\leq (n/b)\cdot 2^{t+2}
$$ 
Since both $i_1$ and $i_2$ are isolated under $\pi$, this means that 
$$
\left|\left\{i\in S\setminus S^*_H\setminus \{j\} \text{~s.t.~}||\pi(j)-\pi(i)||_\infty\leq (n/b)\cdot (2^{t+1}-1)\right\}\right|\leq \GI\cdot \alpha^{d/2} 2^{(t+3)d}\cdot 2^{t+2}+1,
$$
where we used the bound from Definition~\ref{def:isolated} for $i$, but counted the point $i$ itself (this is what makes the bound on $\kappa_j$ weaker than the bound on $\eta_j$). A similar calculation to the one above for $A_1$ now gives
\begin{equation*}
\begin{split}
\kappa_j:=&\sum_{t\geq 0} \max_{\substack{||\pi(j)-\pi(i)||_\infty\geq \\(n/b)\cdot (2^t-1)}} G_{o_i(j)}\cdot \left|\left\{i\in S\setminus S^*_H\setminus \{j\} \text{~s.t.~}||\pi(j)-\pi(i)||_\infty\leq (n/b)\cdot (2^{t+1}-1)\right\}\right|\\
&\leq \sum_{t\geq 0} (\GI\cdot  \alpha^{d/2} 2^{(t+3)d}\cdot 2^{t+2}+1) \min\{1,  2^{-(t-3)\fc}\} \\
&\leq 2^{O(d)}(\GI\cdot \alpha^{d/2}+1)= 2^{O(d)}.
\end{split}
\end{equation*}
We thus have
\begin{equation*}
\begin{split}
A_2\leq &\sum_{j\in \nsq} |y_j|  \kappa_j\leq 2^{O(d)} ||y||_1.
\end{split}
\end{equation*}
Plugging our bounds on $A_1$ and $A_2$ into \eqref{eq:asdf}, we get
\begin{equation*}
\begin{split}
e^{head}_{i}(H, x, \chi)&\leq |G_{o_i(i)}^{-1}|\cdot (A_1+A_2)\leq |G_{o_i(i)}^{-1}|(2^{O(d)}\GI\cdot \alpha^{d/2} ||x'_S||_1+2^{O(d)}||y||_1)\\
&\leq 2^{O(d)}\alpha^{d/2} ||x'_S||_1+\GS \cdot 2^{O(d)} ||y||_1\\
\end{split}
\end{equation*}
as required.

\if 0
Now to obtain the $\ell_\infty$ bound for the case when $\chi_{\nsq\setminus S}=0$, we write for each $i\in S\setminus S^*_H$
\begin{equation}\label{eq:odfjgh}
\begin{split}
A_1(i)\leq &\sum_{t\geq 0} \left(\frac2{1+2^t}\right)^{\fc} \sum_{i'\in S\setminus \{i\}, ||h(i')-h(i)||_\infty\leq 2^t} |x'_{i'}|\\
&\leq ||x'||_\infty\cdot \sum_{t\geq 0} \left(\frac2{1+2^t}\right)^{\fc} \sum_{i\in S\setminus S^*_H}(|S^\pi\cap \B^\infty_{(n/b)\cdot h(i)}((n/b)\cdot 2^{t+2})|-1)\\
\end{split}
\end{equation}
and use the fact that 
$$
|\B^\infty_{h(i)}((n/b)\cdot 2^{t+1})\cap S^\pi|-1\leq \GI\cdot \alpha^{d/2} 2^{(t+3)d}\cdot 2^{t}
$$
for every $i\in (S\setminus S^*_H)$ by definition of an isolated element (see Definition~\ref{def:isolated}). Continuing \eqref{eq:odfjgh}, we get
\begin{equation*}
\begin{split}
A_1\leq ||x'||_\infty\cdot \sum_{t\geq 0} \left(\frac2{1+2^t}\right)^{\fc}\GI\cdot \alpha^{d/2} 2^{(t+3)d}\cdot 2^{t}=||x'||_\infty\cdot \GI 2^{O(d)} \alpha^{d/2},\\
\end{split}
\end{equation*}
implying that 
$$
||e^{head}_{S\setminus S^*_H}(H, x, \chi)||_\infty\leq |G_{o_i(i)}^{-1}|\cdot  A_1\leq 2^{O(d)}\alpha^{d/2} ||x'_S||_\infty
$$
as required.
\fi 
\end{proof}
\begin{remark}
The second bound of this lemma will be useful later in section~\ref{sec:linf} for analyzing \textsc{ReduceInfNorm}.
\end{remark}


We now bound the final error induced by head elements, i.e. $e^{head}(\{H_r\}, x, \chi)$:
\begin{lemma}\label{lm:eh-s-sstar}
Let $x, \chi\in \nsq$, $x'=x-\chi$. \defsk Let $\{\pi_r\}_{r=1}^{r_{max}}$ be a set of permutations, let $H_r=(\pi_r, B, F), F\geq 2d$ be a hashing into $B$ buckets and filter $G$ with sharpness $F$. Let $S^*$ denote the set of elements $i\in S$ that are not isolated under at least $\sqrt{\alpha}$  fraction of $H_r$. Then, one has for $e^{head}$ defined with respect to $S$, 
$$
||e^{head}_{S\setminus S^*}(\{H_r\}, x, \chi)||_1\leq 2^{O(d)}\alpha^{d/2} ||x'_S||_1+\GS \cdot 2^{O(d)} ||\chi_{\nsq\setminus S}||_1.
$$
Furthermore, if $\chi_{\nsq\setminus S}=0$, then $||e^{head}_{S\setminus S^*}(\{H_r\}, x, \chi)||_\infty\leq 2^{d/2}\alpha^{d/2} ||x'_S||_\infty$.
\end{lemma}
\begin{proof}
Recall that by \eqref{eq:eh} one has for each $i\in \nsq$ $e^{head}_i(\{H_r\}, x, \chi)=\quant^{1/5}_{r\in [1:r_{max}]} e^{head}_i(H_r, x, \chi)$. This means that for each $i\in S\setminus S^*$ there exist at least $(1/5-\sqrt{\alpha})r_{max}$ values of $r$ such that 
$e^{head}_i(H_r, x, \chi)>e^{head}_i(\{H_r\}, x, \chi)$, and hence
$$
||e^{head}_{S\setminus S^*}(\{H_r\}, x, \chi)||_1\leq \frac1{(1/5-\sqrt{\alpha})r_{max}}\sum_{r=1}^{r_{max}} ||e^{head}_{S\setminus S^*_r}(H_r, x, \chi)||_1.
$$

By Lemma~\ref{lm:l1b-single-pi} one has
$$
||e^{head}_{S\setminus S^*_{H_r}}(H_r, x, \chi)||_1\leq 2^{O(d)}\alpha^{d/2} ||x'_S||_1+\GS \cdot 2^{O(d)} ||\chi_{\nsq\setminus S}||_1
$$
for all $r$, implying that 
\begin{equation*}
\begin{split}
||e^{head}_{S\setminus S^*}(\{H_r\}, x, \chi)||_1&\leq \frac{1}{(1/5-\sqrt{\alpha})}(2^{O(d)}\alpha^{d/2} ||x'_S||_1+\GS \cdot 2^{O(d)} ||\chi_{\nsq\setminus S}||_1)\\
&\leq 2^{O(d)}\alpha^{d/2} ||x'_S||_1+\GS \cdot 2^{O(d)} ||\chi_{\nsq\setminus S}||_1
\end{split}
\end{equation*}
as required.

The proof of the second bound follows analogously using the $\ell_\infty$ bound from Lemma~\ref{lm:l1b-single-pi}.
\end{proof}
\begin{remark}
The second bound of this lemma will be useful later in section~\ref{sec:linf} for analyzing \textsc{ReduceInfNorm}.
\end{remark}


\subsection{Bounding effect of tail noise}\label{sec:tail-noise}


\begin{lemma}\label{lm:loc-tail-small-single-H}
For any constant $C'>0$ there exists an absolute constant $C>0$ such that for any $x\in \C^{\nsq}$, any integer $k\geq 1$ and $S\subseteq \nsq$ such that $||x_{\nsq\setminus S}||_\infty\leq C'||x_{\nsq\setminus [k]}||_2/\sqrt{k}$, for any integer $B\geq 1$ a power of $2^d$ the following conditions hold.
If $(H, \A)$ are random measurements as in Algorithm~\ref{alg:main-sublinear}, $H=(\pi, B, F)$ satisfies $F\geq 2d$ and $||x_{\nsq\setminus [k]}||_2\geq N^{-\Omega(c)}$, where $O(c)$ is the word precision of our semi-equispaced Fourier transform computation, then for any $i\in \nsq$ one has, for $e^{tail}$ defined with respect to $S$, 
$$
\expect_{H, \A}\left[e^{tail}_{i}(H, \A, x)\right]\leq  \GS \cdot C^d(40+|\H|2^{-\Omega(|\A|)})  ||x_{\nsq \setminus [k]}||_2/\sqrt{B}.
$$ 
\end{lemma}
\begin{proof}

Recall that for any $H=(\pi, B, G), a, \h$ one has $(e^{tail}(H, a\star (\one, \h), x_{\nsq\setminus [k]}))^2=|u_i|^2$, where 
$$
u =\Call{HashToBins}{\widehat{x_{\nsq\setminus S}}, 0, (H, a\star (\one, \h))}.
$$

Since the elements of $\A$ are selected uniformly at random, we have for any $H$ and $\h$ by Lemma~\ref{lm:hashing}, {\bf (3)}, since $a\star (\one, \h)$ is uniformly random in $\nsq$, that
\begin{equation}\label{eq:expect-a-0hg443g}
\expect_a[(e^{tail}_i(H, a\star (\one, \h), x))^2]=\expect_{a}[\abs{G_{o_i(i)}^{-1}\omega^{-(a\star (\one, \h))^T\Sigma i}u_{h(i)} - x_i}^2]\leq \mu^2_{H, i}(x)+N^{-\Omega(c)},
\end{equation}
where $c>0$ is the large constant that governs the precision of our Fourier transform computations. 
By Lemma~\ref{lm:hashing}, {\bf (2)} applied to the pair $(\widehat{x_{\nsq\setminus S}}, 0)$ there exists a constant $C>0$ such that
$$
\expect_H[\mu^2_{H, i}] \leq  \GSS\cdot C^d ||x_{\nsq \setminus S}||_2^2/B
$$ 
We would like to upper bound the rhs in terms of $||x_{\nsq\setminus [k]}||_2^2$ (the tail energy), but this requires an argument since $S$ is not exactly the set of top $k$ elements of $x$. However, since $S$ contains the large coefficients of $x$, a bound is easy to obtain. Indeed, denoting the set of top $k$ coefficients of $x$ by $[k]\subseteq \nsq$ as usual, we get
\begin{equation*}
||x_{\nsq \setminus S}||_2^2\leq ||x_{\nsq \setminus (S\cup [k])}||_2^2+||x_{[k]\setminus S}||_2^2\leq ||x_{\nsq \setminus [k]}||_2^2+k\cdot ||x_{[k]\setminus S}||_\infty^2\leq (C'+1)||x_{\nsq \setminus [k]}||_2.
\end{equation*}
Thus, we have  
\begin{equation*}
\expect_H[\mu^2_{H, i}(x)+N^{-\Omega(c)}]\leq \GSS\cdot (C'+2)C^d ||x_{\nsq \setminus [k]}||_2^2/B,
\end{equation*}
where we used the assumption that $||x_{\nsq \setminus k}||_2\geq N^{-\Omega(c)}$.
We now get by Jensen's inequality 
\begin{equation}\label{eq:mu-b}
\expect_H[\mu_{H, i}(x)]\leq\GS\cdot (C'')^d ||x_{\nsq \setminus k}||_2/\sqrt{B}
\end{equation}
for a constant $C''>0$. Note that 


By ~\eqref{eq:expect-a-0hg443g} for each $i\in \nsq$, hashing $H$, evaluation point $a\in \nsq\times \nsq$ and direction $\h$ we have
$\expect_a[(e^{tail}_i(H, a\star (\one,  \h), x))^2]=(\mu_{H, i}(x))^2$. Applying Jensen's inequality, we hence get for any $H$ and $\h\in \H$
\begin{equation}
\expect_a[e^{tail}_{i}(H, a\star (\one, \h), x)]\leq \mu_{H, i}(x).
\end{equation}
 Applying Lemma~\ref{lm:quant-exp} with $Y=e^{tail}_{i}(H, a\star (\one, \h), x)$ and $\gamma=1/5$ (recall that the definition of $e^{tail}_i(H, z, x)$ involves a $1/5$-quantile over $\A$) and using the previous bound, we get, for any fixed $H$ and $\h\in \H$
\begin{equation}\label{eq:e-b-exp}
\expect_\A\left[\left|e^{tail}_{i}(H, \A\star (\one, \h), x)-40\cdot \mu_{H, i}(x)\right|_+\right]\leq \mu_{H, i}(x)\cdot 2^{-\Omega(|\A|)},
\end{equation}
and hence by a union bound over all $\h\in \H$ we have 
\begin{equation*}
\expect_\A\left[\sum_{\h\in \H}\left|e^{tail}_{i}(H, \A\star (\one, \h), x)-40\cdot \mu_{H, i}(x)\right|_+\right]\leq \mu_{H, i}(x)\cdot |\H|2^{-\Omega(|\A|)}.
\end{equation*}
Putting this together with~\eqref{eq:mu-b}, we get 
\begin{equation*}
\begin{split}
&\expect_{H, \A}\left[e^{tail}_i(H, \A, x)\right]\\
&=\expect_H\left[\expect_\A\left[40\mu_{H, i}(x)+\sum_{\h\in \H}\left|e^{tail}_{i}(H, \A\star (\one, \h), x)-40\cdot \mu_{H, i}(x)\right|_+\right]\right]\\
&\leq\expect_H\left[\mu_{H, i}(x)(40+|\H|2^{-\Omega(|\A|)})\right]\\
&\leq \GS (C'')^d (40+|\H|2^{-\Omega(|\A|)}) ||x_{\nsq \setminus k}||_2/\sqrt{B}
\end{split}
\end{equation*}
as required.
\end{proof}


\begin{lemma}\label{lm:loc-tail-small}
For any constant $C'>0$ there exists an absolute constant $C>0$ such that for any $x\in \C^{\nsq}$, any integer $k\geq 1$ and $S\subseteq \nsq$ such that $||x_{\nsq\setminus S}||_\infty\leq C'||x_{\nsq\setminus [k]}||/\sqrt{k}$, if $B\geq 1$, then the following conditions hold , for $e^{tail}$ defined with respect to $S$.


If hashings $H_r=(\pi_r, B, F), F\geq 2d$ and sets $\A_r, |\A_r|\geq c_{max}$ for $r=1,\ldots, r_{max}$ are chosen at random, then
\begin{description}
\item[(1)] for every $i\in \nsq$ one has 
$$
\expect_{\{(H_r, \A_r)\}}\left[e^{tail}_i(\{H_r, \A_r\}, x)\right]\leq  \GS C^d (40+|\H|2^{-\Omega(c_{max})})  ||x_{\nsq \setminus [k]}||_2/\sqrt{B}.
$$ 

\item[(2)] for every $i\in \nsq$ one has 
$$
\prob_{\{(H_r, \A_r)\}}\left[e^{tail}_i(\{H_r, \A_r\}, x)>  \GS C^d (40+|\H|2^{-\Omega(c_{max})})  ||x_{\nsq \setminus [k]}||_2/\sqrt{B}\right]=2^{-\Omega(r_{max})}
$$ 
and
\begin{equation*}
\begin{split}
&\expect_{\{(H_r, \A_r)\}}\left[\left|e^{tail}_i(\{H_r, \A_r\}, x)- \GS C^d (40+|\H|2^{-\Omega(c_{max})})  ||x_{\nsq \setminus [k]}||_2/\sqrt{B}\right|_+\right]\\
&=2^{-\Omega(r_{max})}\cdot \GS C^d (40+|\H|2^{-\Omega(c_{max})})  ||x_{\nsq \setminus [k]}||_2/\sqrt{B}.
\end{split}
\end{equation*}

\end{description}
\end{lemma}
\begin{proof}

Follows by applying  Lemma~\ref{lm:quant-exp} with $Y=e^{tail}_{i}(H_r, \A_r, x)$.
\end{proof}

\subsection{Putting it together}
The bounds from the previous two sections yield a proof of Theorem~\ref{thm:l1-res-loc}, which we restate here for convenience of the reader:

{\em {\bf Theorem~\ref{thm:l1-res-loc}}
For any constant $C'>0$ there exist absolute constants $C_1, C_2, C_3>0$ such that for any $x\in \C^\nsq$, any integer $k\geq 1$ and any $S\subseteq \nsq$ such that $||x_{\nsq\setminus S}||_\infty\leq C'\mu$, where $\mu=||x_{\nsq\setminus [k]}||_2/\sqrt{k}$, the following conditions hold.

Let $\pi_r=(\Sigma_r, q_r), r=1,\ldots, r_{max}$ denote permutations, and let $H_r=(\pi_r, B, F)$, $F\geq 2d$, where $B\geq \GT k/\alpha^d$ for $\alpha\in (0, 1)$ smaller than a constant. Let $S^*\subseteq S$ denote the set of elements that are not isolated with respect to at least a $\sqrt{\alpha}$ fraction of hashings $\{H_r\}$. Then if $r_{max}, c_{max}\geq (C_1/\sqrt{\alpha})\log\log N$, then with probability at least $1-1/\log^2 N$ over the randomness of the measurements for all $\chi\in \C^\nsq$ such that $x':=x-\chi$ satisifies $||x'||_\infty/\mu\leq N^{O(1)}$ one has
$$
L:=\bigcup_{r=1}^{r_{max}}\textsc{LocateSignal}\left(\chi, k, \{m(\wh{x}, H_r, a\star (\one, \h))\}_{r=1, a\in \A_r, \h\in \H}^{r_{max}}\right)
$$
satisfies
$$
||x'_{S\setminus S^*\setminus L}||_1\leq (C_2\alpha)^{d/2} ||x'_S||_1+C_3^{d^2}(||\chi_{\nsq\setminus S}||_1+||x'_{S^*}||_1)+4\mu |S|.
$$

}
\begin{proof}
First note that with probability at least  $1-1/(10\log^2 N)$  for every $s\in [1:d]$ the sets $\A_r\star (\zero, \mathbf{e}_s)$ are balanced (as per Definition~\ref{def:balance}) for all $r=1,\ldots, r_{max}$ and all $\h\in \H$ by Claim~\ref{cl:balanced}.

By Corollary~\ref{cor:loc} applied with $S'=S\setminus S^*$ one has 
$$
||(x-\chi)_{(S\setminus S^*)\setminus L}||_1\leq 20 \cdot (||e^{head}_{S\setminus S^*}(\{H_r\}, x')||_1+||e^{tail}(\{H_r, \A_r\}, x)||_1)+||x'||_\infty |S|\cdot N^{-\Omega(c)}.
$$

We also have  
$$
||e^{head}_{S\setminus S^*}(\{H_r\}, x')||_1\leq 2^{O(d)}\alpha^{d/2} ||x'_S||_1+\GS \cdot 2^{O(d)} ||\chi_{\nsq\setminus S}||_1
$$
by Lemma~\ref{lm:eh-s-sstar} and with probability at least $1-1/(10\log^2 N)$
\begin{equation*}
||e^{tail}_{S\setminus S^*}(\{H_r, \A_r\}, x)||_1\leq \GS C^d (40+|\H|2^{-\Omega(c_{max})})  ||x_{\nsq \setminus [k]}||_2|S|/\sqrt{B} 
\end{equation*}
by Lemma~\ref{lm:loc-tail-small}. The rhs of the previous equation is bounded by $|S|\mu$ by the choice of $B$ as long as $\alpha$ is smaller than a absolute constant, as required.
Putting these bounds together and using the fact that $|\H|\leq \log N$ (so that $|\H|\cdot (2^{-\Omega(r_{max})}+2^{-\Omega(c_{max})})\leq 1$), and taking a union bound over the failure events, we get the result.
\end{proof}

